\documentclass[10pt,letterpaper]{article}
\usepackage[margin=0.85in]{geometry}
\usepackage{amsmath,amssymb,amsthm}
\usepackage{booktabs}
\usepackage{multirow}
\usepackage{graphicx}
\usepackage{caption}
\captionsetup{font=small,labelfont=bf}
\setcounter{secnumdepth}{2}

\newtheorem*{slemma}{Lemma}
\newtheorem*{sproposition}{Proposition}
\newtheorem*{stheorem}{Theorem}
\newtheorem*{sdefinition}{Definition}

\title{Supplementary Material for ``Affordable Real Style Transfer:\\
Training Spectral Flow Matching on an RTX 3060 in Minutes''}
\author{Anonymous Submission}
\date{}

\begin{document}
\maketitle

\appendix

\section{Method Details}
\label{sec:method_details}

This section expands Section~3 of the main paper. All theorem, lemma, proposition, and definition numbers refer to the main paper.

\subsection{Problem Setup}
\label{sec:problem}

Let $x_c \in \mathbb{R}^{H \times W \times 3}$ be a content image and $s \in \{1,\ldots,K\}$ a target style domain. We learn a mapping $G_\theta(x_c, s) = \mathcal{D}(\hat z_s)$, where $\mathcal{D}$ is a frozen VAE decoder and $\hat z_s$ is a stylized latent. The transport is a deterministic ODE (rectified flow) from $z_0 = z_c$ to $z_1 \sim p_{\text{style}}$.

\paragraph{VAE latent choice.}
We work in the latent space of a pre-trained SDXL VAE, $z_c = \mathcal{E}(x_c) \in \mathbb{R}^{4 \times 32 \times 32}$. Three properties motivate this choice. First, the decoder is frozen and perceptually calibrated, so a small edit in latent space maps to a semantically meaningful edit in pixel space. Second, the $32 \times 32$ spatial size yields $16 \times 16$ subbands after a single-level Haar DWT, which a small backbone can process at batch~24 on a 12\,GB GPU. Third, four channels match the standard SDXL latent convention, allowing reuse of off-the-shelf VAE weights without retraining.

\subsection{Haar DWT and Spectral Decoupling}
\label{sec:haar}

We restate Lemma~1 and Proposition~1 of the main paper and give full proofs.

\begin{slemma}[Wavelet isometry and Parseval identity --- Lemma~1 of the main paper]
The single-level 2D Haar DWT $\mathcal{W}: \mathbb{R}^{C\times H\times W} \to \mathbb{R}^{C\times H/2\times W/2}\!\times\!\{4\text{ subbands}\}$ is an isometry: $\|\mathcal{W}(z)\|_F^2 = \|z\|_F^2$. For any partition of the four subbands into disjoint sets $A,B$, $\|z\|_F^2 = \|A\|_F^2 + \|B\|_F^2$ with no cross term.
\end{slemma}
\begin{proof}
The Haar kernel is $H_2 = \tfrac{1}{\sqrt{2}}\bigl[\begin{smallmatrix}1 & 1\\ 1 & -1\end{smallmatrix}\bigr]$, satisfying $H_2^\top H_2 = I_2$. The 2D transform applies $H_2$ separably along rows then columns, so the per-patch transform is the Kronecker product $H_2 \otimes H_2$. Kronecker products preserve orthogonality: $(A\otimes B)^\top(A\otimes B) = (A^\top A)\otimes(B^\top B) = I_2 \otimes I_2 = I_4$. Hence $\mathcal{W}$ is an orthogonal linear map and $\|\mathcal{W}(z)\|_F^2 = \|z\|_F^2$ (Parseval). Orthogonality of the four subband blocks gives $\|z\|_F^2 = \|\text{LL}\|_F^2 + \|\text{LH}\|_F^2 + \|\text{HL}\|_F^2 + \|\text{HH}\|_F^2$, so any partition $A \cup B$ satisfies $\|z\|_F^2 = \|A\|_F^2 + \|B\|_F^2$ with no cross term.
\end{proof}

\begin{sproposition}[Spectral content--style decoupling --- Proposition~1 of the main paper]
If the content signal concentrates its energy in LL, $\|z^{\text{cnt}}_{\text{LL}}\|_F^2 \geq (1-\epsilon)\|z^{\text{cnt}}\|_F^2$, then any transform $T$ acting only on $\{\text{LH,HL,HH}\}$ perturbs the content energy by at most $\epsilon\,\|z^{\text{cnt}}\|_F^2$.
\end{sproposition}
\begin{proof}
By Lemma~1, the LL and HF subspaces are orthogonal, so $\|T(z)-z\|_F^2 = \|T(z^{\text{sty}})-z^{\text{sty}}\|_F^2 + \|T(z^{\text{cnt}}_{\text{HF}}) - z^{\text{cnt}}_{\text{HF}}\|_F^2$. Since $T$ acts only on HF, $T(z^{\text{cnt}}_{\text{LL}}) = z^{\text{cnt}}_{\text{LL}}$, and the second term is bounded by $\|z^{\text{cnt}}_{\text{HF}}\|_F^2 \leq \epsilon\,\|z^{\text{cnt}}\|_F^2$ by the energy-concentration assumption.
\end{proof}

\subsection{Base Locking}
\label{sec:base_locking}

We set the LL velocity loss weight $w_{\text{LL}}=0$ during training but still apply the small predicted LL drift at inference. The training loss is
\[
\mathcal{L}(\theta) = \sum_{c \in \{\text{LH, HL}\}} w_c \, \mathbb{E}_{t} \big\| v_\theta^{c}(z_t, t, s) - u_t^{c} \big\|_2^2,
\qquad w_{\text{LH}}=w_{\text{HL}}=1,\; w_{\text{HH}}=0.
\]

\paragraph{Why $w_{\text{LL}}=0$.}
By Proposition~1, content energy concentrates in LL. Training an LL velocity would couple content-reconstruction gradients with style gradients in the same subspace, re-creating the degeneration attractor (Theorem~1) inside the LL block: the LL head would be pulled toward the style-averaged velocity $\bar v_\theta$, the only critical point at which content and style gradients cancel. Setting $w_{\text{LL}}=0$ removes the LL term from the flow-matching loss, so no gradient flows to $v_\theta^{\text{LL}}$ through $\mathcal{L}_{\text{FM}}$. The LL head drifts only through shared backbone weights and remains small in norm.

\paragraph{Why LL is not frozen at inference.}
LL also carries global style (color palette, lighting). Freezing LL at inference removes the only path by which the model can shift global tone, costing $0.0141$ CLIP-S (Table~\ref{tab:ablation_supp}). Letting the small freely-drifting predicted LL velocity be applied at inference gives both the content anchor (LL is barely trained, so $v_\theta^{\text{LL}}$ is small) and a global-style channel. The two objectives --- content preservation in LL and style injection in HF --- no longer compete because they act on orthogonal subspaces (Lemma~1).

\subsection{Fiber Flow and WCT}
\label{sec:fiber}

\begin{sdefinition}[Spectral fiber bundle --- Definition~1 of the main paper]
The total space $E = \mathbb{R}^{C\times H\times W}$ is equipped with the projection $\pi: E \to B$ that returns the LL subband, $\pi(z) = z^{\text{LL}} \in B$. The fiber above $b \in B$ is $F_b = \{z : z^{\text{LL}}=b\}$, parameterized by $(\text{LH,HL,HH})$. A map $T:E\to E$ is \emph{fiber-preserving} if $\pi(T(z)) = \pi(z)$ for all $z$.
\end{sdefinition}

\begin{sproposition}[WCT is fiber-preserving --- Proposition~2 of the main paper]
Per-subband whitening-and-coloring on each high-frequency subband leaves LL unchanged, hence $\pi(\mathrm{WCT}(\hat z_s)) = \pi(\hat z_s) = \hat\ell$. The content anchor is an invariant of style injection; style transfer is a vertical map along the fibers $F_b$.
\end{sproposition}
\begin{proof}
For each $c \in \{\text{lh,hl,hh}\}$, $\mathrm{WCT}_c(\hat c, c^{\text{ref}}) = \Sigma_{c^{\text{ref}}}^{1/2}\, \Sigma_{\hat c}^{-1/2}\, (\hat c - \mu_{\hat c}) + \mu_{c^{\text{ref}}}$ depends only on the HF subband $c$. LL is not an argument of any $\mathrm{WCT}_c$, so $\ell^{\text{new}} = \hat\ell$ by construction. Therefore $\pi(\mathrm{WCT}(\hat z_s)) = \hat\ell = \pi(\hat z_s)$, which is the fiber-preserving condition. By Lemma~1, the four subbands are orthogonal, so per-fiber statistics matching is equivalent to per-frequency-band style injection in the original latent, with no cross-band leakage.
\end{proof}

\subsection{End-of-Trajectory Injection}
\label{sec:eota}

\begin{sproposition}[$\alpha$-identifiability under endpoint injection --- Proposition~3 of the main paper]
Let $r_n(\alpha) = (1-\alpha)^n$ be the content residual after $n$ per-step injections with blend $\alpha\in(0,1]$. Under per-step injection with $n\geq 2$, the set of $\delta$-indistinguishable pairs covers $[0,1]^2$ for any $\delta > (1-\alpha_{\min})^n$. Under endpoint injection ($n=1$), $r_1(\alpha)=1-\alpha$ is injective, so all pairs are distinguishable for any $\delta < 1$.
\end{sproposition}
\begin{proof}
For $n \geq 2$, $r_n(\alpha) = (1-\alpha)^n$ is convex and decreasing on $[0,1]$ with $r_n(0)=1$, $r_n(1)=0$, and second derivative $n(n-1)(1-\alpha)^{n-2} \geq 0$. Its derivative magnitude is $|r_n'(\alpha)| = n(1-\alpha)^{n-1}$, which tends to $0$ as $\alpha \to 1$: near $\alpha=1$, $r_n$ is flat, so a continuum of $\alpha$ values produce residuals within $\delta$ of each other. On $[\alpha_{\min},1]$ the Lipschitz constant is $n(1-\alpha_{\min})^{n-1}$; for $n=8$, $\alpha_{\min}=0.5$ this is $8 \cdot 0.5^7 \approx 0.0625$, so resolution below $0.0625$ is lost. Hence for any $\delta > (1-\alpha_{\min})^n$ the $\delta$-indistinguishable pairs cover the square. For $n=1$, $r_1(\alpha) = 1-\alpha$ is affine with slope $-1$, hence strictly monotone and injective. Distinct $\alpha$ therefore yield distinct residuals, and any two blends are distinguishable for every $\delta < 1$.
\end{proof}

This is why style is injected only at the last ODE step: it restores $\alpha$ as a meaningful trade-off knob and matches the Schr\"odinger-bridge view in which style is a terminal constraint, not a per-step perturbation.

\subsection{Stochastic Frequency Routing}
\label{sec:routing}

The cross-attention that routes style tokens into the velocity field can operate on the full latent (Euclidean query) or on the high-frequency tokens only (wavelet query). At inference we always use the wavelet query. At training we draw $B \sim \mathrm{Bernoulli}(p)$ per step and use the wavelet query with probability $p$. The sweep $p \in \{0.0, 0.5, 0.8, 1.0\}$ gives
\[
\text{CLIP-S} \in \{0.7061,\, 0.7083,\, 0.7213,\, 0.7226\},\quad
\text{LPIPS} \in \{0.2606,\, 0.2480,\, 0.2868,\, 0.3068\}.
\]

\paragraph{Why $p=0.8$.}
Two failure modes bracket the sweep. At $p=0.0$ the model never sees a wavelet query during training, so the query projection is mis-specified at inference and CLIP-S drops to $0.7061$ (a train/test mismatch). At $p=1.0$ the style encoder over-specializes to high-frequency statistics and loses global-palette information, which inflates LPIPS to $0.3068$ through content tearing. The point $p=0.8$ is the unique sweep value that clears both thresholds (CLIP-S~$>~0.72$ and LPIPS~$<~0.30$). Mixing the two query modes regularizes the projection: biasing toward the inference-time (wavelet) query while retaining enough Euclidean exposure to preserve palette information.

\subsection{Degeneration Attractor}
\label{sec:attractor_supp}

\begin{sdefinition}[Degeneration manifold --- Definition~2 of the main paper]
$\mathcal{M}_0 = \big\{\theta \,:\, v_\theta(z, t, s) = \bar v_\theta(z, t) = \mathbb{E}_{s}[v_\theta(z, t, s)] \;\;\forall s, z, t\big\}$, equivalently $\nabla_s v_\theta|_{\theta\in\mathcal{M}_0} = 0$.
\end{sdefinition}

\begin{stheorem}[$\mathcal{M}_0$ is a local minimum under multiplicative attenuation --- Theorem~1 of the main paper]
Assume $\mathcal{L} = w_{\text{FM}} \mathcal{L}_{\text{FM}} + w_{\text{SWD}} \mathcal{L}_{\text{SWD}}$ with $w_{\text{FM}} \gg w_{\text{SWD}}$, and three multiplicative attenuation factors act on the style gradient: a scalar gate $g(\theta)\in[0,1]$, an attention-entropy factor $h(\theta)\in[0,1]$ ($h=1$ uniform), and a normalization washout $n(\theta)\in[0,1]$. Let $\kappa(\theta) = g(\theta)\,h(\theta)\,n(\theta)$. If at $\theta_0\in\mathcal{M}_0$: (i) $\kappa(\theta_0) \leq \sqrt{w_{\text{SWD}}/w_{\text{FM}}}$; (ii) $\nabla^2_{\theta_\perp} \mathcal{L}_{\text{FM}}|_{\theta_0}$ is positive definite on $T_{\theta_0}^\perp \mathcal{M}_0$; (iii) $\|\nabla^2 \mathcal{L}_{\text{SWD}}|_{\theta_0}\|_{\text{op}} \leq \kappa(\theta_0)\,\|\nabla^2 \mathcal{L}_{\text{FM}}\|_{\text{op}}$; then $\theta_0$ is a strict local minimum of $\mathcal{L}$, with basin of attraction of radius $\Theta(\kappa(\theta_0))$.
\end{stheorem}
\begin{proof}
\emph{Gradient attenuation.} The scalar gate $g$ multiplies the style-conditioned signal in the velocity field; uniform attention (entropy factor $h\!\to\!1$) spreads style tokens evenly so their weighted sum washes out the style direction; GroupNorm (factor $n$) removes per-channel mean and variance, which are precisely the channels that carry style statistics. Composing these, the style gradient satisfies $\|\nabla_\theta \mathcal{L}_{\text{SWD}}\| \leq g\,h\,n \,\|\nabla_\theta \mathcal{L}_{\text{FM}}\| = \kappa\,\|\nabla_\theta \mathcal{L}_{\text{FM}}\|$.

\emph{First-order condition.} On $\mathcal{M}_0$, $v_\theta = \bar v_\theta$ is style-independent, i.e.\ the conditional mean over styles. This is the minimizer of $\mathcal{L}_{\text{FM}}$ averaged over $s$, so $\nabla_\theta \mathcal{L}_{\text{FM}}|_{\theta_0} = 0$. By condition (i), $\|\nabla_\theta \mathcal{L}_{\text{SWD}}\| \leq \kappa\,\|\nabla_\theta \mathcal{L}_{\text{FM}}\| = 0$, so $\nabla \mathcal{L}(\theta_0) = w_{\text{FM}}\nabla\mathcal{L}_{\text{FM}} + w_{\text{SWD}}\nabla\mathcal{L}_{\text{SWD}} = 0$.

\emph{Second-order condition.} Decompose a perturbation as $\theta = \theta_0 + \delta\theta_\perp + \delta\theta_\parallel$ with $\delta\theta_\perp \in T^\perp_{\theta_0}\mathcal{M}_0$. The second-order expansion along the normal direction is
\[
\delta^2 \mathcal{L} = \tfrac12\,\delta\theta_\perp^\top \bigl(\nabla^2_\perp \mathcal{L}_{\text{FM}} + (w_{\text{SWD}}/w_{\text{FM}})\nabla^2_\perp \mathcal{L}_{\text{SWD}}\bigr)\delta\theta_\perp.
\]
By (ii), $\lambda_{\min}(\nabla^2_\perp \mathcal{L}_{\text{FM}}) > 0$. By (iii), $\|\nabla^2_\perp \mathcal{L}_{\text{SWD}}\|_{\text{op}} \leq \kappa\,\|\nabla^2 \mathcal{L}_{\text{FM}}\|_{\text{op}}$. Therefore
\[
\delta^2 \mathcal{L} \geq \tfrac12 \bigl(\lambda_{\min}(\nabla^2_\perp \mathcal{L}_{\text{FM}}) - (w_{\text{SWD}}/w_{\text{FM}})\,\kappa\,\|\nabla^2 \mathcal{L}_{\text{FM}}\|_{\text{op}}\bigr)\|\delta\theta_\perp\|^2.
\]
By (i), $(w_{\text{SWD}}/w_{\text{FM}})\,\kappa \leq \kappa^{2}/\kappa \leq \kappa$, so the bracket is positive whenever $\kappa < \lambda_{\min}/\|\nabla^2 \mathcal{L}_{\text{FM}}\|_{\text{op}}$, which holds within radius $O(\kappa)$ of $\theta_0$. Hence $\delta^2\mathcal{L} > 0$ on $T^\perp\mathcal{M}_0$ in a neighborhood of $\theta_0$, making $\theta_0$ a strict local minimum. The basin radius scales as $\Theta(\kappa)$.
\end{proof}

\paragraph{Empirical attenuation.}
In the Euclidean baseline the three factors are measured as $g \approx 0.05$, $h \approx 1$ (uniform attention), $n \approx 0.32$ (GroupNorm washout), giving $\kappa \approx 0.016$. With $w_{\text{SWD}}/w_{\text{FM}} = 0.01$, $\sqrt{w_{\text{SWD}}/w_{\text{FM}}} = 0.1 \gg \kappa$, so condition (i) holds and the basin radius is $\Theta(0.016)$: the model is trapped. Because the attenuation is multiplicative, removing any single mechanism (e.g.\ re-initializing the gate) is cancelled by the others.

\begin{sproposition}[Wavelet escape --- Proposition~4 of the main paper]
Under the spectral decomposition, $\nabla_\theta \mathcal{L}_{\text{FM}}$ flows through LL and $\nabla_\theta \mathcal{L}_{\text{SWD}}$ flows through LH/HL/HH. By Proposition~1 these subspaces are orthogonal, so $\nabla^2\mathcal{L}$ block-diagonalizes and the cross term $\nabla^2_{\text{LL, HF}}\mathcal{L}$ vanishes. Condition (ii) of Theorem~1 no longer binds the HF block, and $\mathcal{M}_0$ ceases to be a critical manifold of $\mathcal{L}_{\text{wavelet}}$.
\end{sproposition}

\subsection{Architecture and Training Hyperparameters}
\label{sec:hyper}

Table~\ref{tab:hyper} lists all hyperparameters of the SFM baseline.

\begin{table}[h]
\centering
\caption{SFM architecture and training hyperparameters.}
\label{tab:hyper}
\small
\begin{tabular}{@{}ll@{}}
\toprule
Parameter & Value \\
\midrule
Latent space & SDXL VAE, $4 \times 32 \times 32$ \\
Wavelet & Haar, 3 levels \\
Backbone & 4 residual blocks, $C=64$, AdaIN conditioning \\
Subband heads & $3\times3$ conv per subband (LH, HL; HH dropped) \\
Concatenated input & $[z^{\text{LL}}, z^{\text{LH}}, z^{\text{HL}}] \in \mathbb{R}^{12 \times 16 \times 16}$ \\
Style encoder & frozen DINOv2 + learnable token bank (5 tokens $\times$ 64 dims) \\
Loss weights & $w_{\text{LH}}=w_{\text{HL}}=1$,\; $w_{\text{LL}}=0$,\; $w_{\text{HH}}=0$ \\
Stochastic routing & Bernoulli, $p=0.8$ \\
ODE solver & Heun (2nd-order Runge--Kutta), 8 steps \\
Style injection & per-subband WCT, endpoint only; $\Sigma + 10^{-5}I$ \\
Optimizer & Adam, learning rate $10^{-4}$ \\
Batch size & 24 \\
Epochs & 5 \\
Hardware & single RTX 3060 (12\,GB) \\
Trainable parameters & 903{,}248 \\
Training time & 3.08 min end-to-end \\
% TODO: confirm the endpoint blend coefficient alpha used for the main-result checkpoint.
\bottomrule
\end{tabular}
\end{table}

\section{Subtractive Ablation Audit}
\label{sec:ablation_supp}

This section documents the ablation audit referenced in the main paper (Section~4.4). Each entry removes or replaces one component of the SFM baseline and reports the resulting 750-pair Distinct5-WikiArt-512 metrics. All numbers are single-seed. $\Delta$ values are relative to the SFM baseline (CLIP-S $0.7213$, LPIPS $0.2868$).

\subsection{Component Ablation Table}
\label{sec:ablation_table}

Table~\ref{tab:ablation_supp} reproduces the main paper's ablation table and adds variants named in the main text.

\begin{table}[h]
\centering
\caption{Subtractive ablation. Each row removes or replaces one component of SFM.}
\label{tab:ablation_supp}
\small
\setlength{\tabcolsep}{4pt}
\begin{tabular}{@{}lcccl@{}}
\toprule
Variant & CLIP-S & LPIPS & $\Delta$CLIP-S & Note \\
\midrule
\textbf{SFM (full)} & \textbf{0.7213} & \textbf{0.2868} & --- & baseline \\
\midrule
\multicolumn{5}{l}{\emph{Wavelet decomposition}} \\
\quad No DWT (Euclidean) & 0.7117 & 0.2994 & $-$0.0096 & attractor active \\
\quad DWT level 1 (vs.\ 3) & 0.7261 & 0.3296 & $+$0.0048 & under-decomposed \\
\quad DWT level 4 (vs.\ 3) & 0.7316 & 0.3461 & $+$0.0103 & over-decomposed, loses position \\
\quad Daubechies-2 (vs.\ Haar) & 0.7258 & 0.3288 & $+$0.0045 & basis function not critical \\
\midrule
\multicolumn{5}{l}{\emph{Base locking}} \\
\quad LL lock at inference & 0.7174 & 0.3372 & $-$0.0039 & kills global style channel \\
\quad $w_{\text{LL}}=1.0$ at train & 0.7174 & 0.2774 & $-$0.0039 & content-heavy trade-off \\
\quad Train HH head & 0.7213 & 0.2868 & $\pm$0.0001 & HH head removed (dead) \\
\midrule
\multicolumn{5}{l}{\emph{Style injection}} \\
\quad Per-step AdaIN (not EOTA) & 0.7361 & 0.3843 & $+$0.0148 & $\alpha$-decay, LPIPS explodes \\
\quad AdaIN (diagonal) vs.\ WCT & 0.7266 & 0.3402 & $+$0.0053 & loses channel correlation \\
\quad Inject LL too & 0.7215 & 0.3856 & $+$0.0002 & color shift propagates \\
\midrule
\multicolumn{5}{l}{\emph{Stochastic routing}} \\
\quad $p=0.0$ (Euclidean train) & 0.7061 & 0.2606 & $-$0.0152 & train/test mismatch \\
\quad $p=0.5$ & 0.7083 & 0.2480 & $-$0.0130 & under-trained DWT route \\
\quad $p=1.0$ (always DWT) & 0.7226 & 0.3068 & $+$0.0013 & style encoder over-specialized \\
\midrule
\multicolumn{5}{l}{\emph{Solver}} \\
\quad Euler (vs.\ Heun) & 0.7210 & 0.3129 & $-$0.0003 & structural DOF lost \\
\quad RK4 (vs.\ Heun) & 0.7265 & 0.3235 & $+$0.0052 & saturated, no gain \\
\midrule
\multicolumn{5}{l}{\emph{Capacity (Pareto-mapping knobs)}} \\
\quad depth $=2$ (vs.\ 4) & 0.7172 & 0.2705 & $-$0.0041 & 1:8 trade-off \\
\quad dim $=32$ (vs.\ 64) & 0.7153 & 0.2705 & $-$0.0060 & 1:8 trade-off \\
\quad gate $=0$ (no gate) & 0.7080 & 0.2641 & $-$0.0133 & 1:8 trade-off \\
\quad depth $=6$ (vs.\ 4) & NaN & NaN & --- & WCT \texttt{eigh} unstable \\
\midrule
\multicolumn{5}{l}{\emph{Added modules (removed with no gain)}} \\
\quad DINO feature extraction (added) & 0.7088 & 0.2868 & $-$0.0125 & redundant in wavelet regime \\
\quad dim $=96$ (added) & --- & --- & --- & underfitting; removed \\
\quad Terminal SWD (added) & 0.7213 & 0.2868 & $\pm$0.0001 & redundant \\
\quad HF loss weighting (added) & 0.7213 & 0.2868 & $\pm$0.0001 & redundant \\
\bottomrule
\end{tabular}
\end{table}

\subsection{Removed Modules List}
\label{sec:removed_modules}

Table~\ref{tab:removed} lists modules that were added to the SFM baseline during development and later removed because they produced no measurable gain or were harmful. The pattern is consistent: any module operating in the Euclidean regime is pulled to the degeneration attractor; any module operating in the wavelet regime is either redundant (the decomposition already disentangles content and style) or harmful (it re-entangles them).

\begin{table}[h]
\centering
\caption{Added modules tested and removed from the SFM baseline. The full 22-entry ledger is released with the code; entries marked TBR are documented in the released ablation log.}
\label{tab:removed}
\small
\begin{tabular}{@{}lll@{}}
\toprule
\# & Module added then removed & Effect when added \\
\midrule
1 & DINO feature extraction & $-$0.0125 CLIP-S \\
2 & Per-step AdaIN (replaces EOTA) & $+$0.015 CLIP-S, $+$0.10 LPIPS \\
3 & HH-head training & $\pm$0.0001 \\
4 & High-frequency loss weighting & $\pm$0.0001 \\
5 & Terminal SWD term & $\pm$0.0001 \\
6 & depth $=6$ backbone & NaN (WCT unstable) \\
7 & dim $=96$ backbone & underfitting \\
\midrule
8--22 & 15 additional variants & no measurable gain \\
      & (alternative style encoders, & (full per-module deltas \\
      & auxiliary losses, multi-scale & released with code) \\
      & fusion, normalization schemes) & \\
\bottomrule
\end{tabular}
% TODO: replace entries 8--22 with the named modules and per-module deltas
% from the released ablation log; the grouping above is a placeholder.
\end{table}

\subsection{Pareto 1:8 Trade-off Analysis}
\label{sec:pareto}

The capacity knobs (depth, width, gate magnitude) all move the operating point along a single line with slope $\approx 8$: reducing capacity improves LPIPS by roughly $8\times$ the CLIP-S cost. The same line is traced by the endpoint blend $\alpha$, the schedule shape, and the training duration. None of these knobs breaks the frontier.

\paragraph{Why capacity maps to one line.}
On $\mathcal{M}_0$ (Definition~2 of the main paper) the velocity field factorizes as $v_\theta(z,t,s) = \bar v(z,t) + \kappa(\theta)\,\delta v_s(z,t)$, where $\kappa = g\,h\,n$ is the multiplicative attenuation of Theorem~1. The content term $\mathcal{L}_{\text{FM}}$ is independent of $\kappa$; the style term $\mathcal{L}_{\text{SWD}}$ scales as $\kappa^2$. Reducing capacity reduces $\|\delta v_s\|$ by a factor $\beta$, which reduces CLIP-S by $O(\beta)$ and reduces LPIPS by $O(\beta)$ as well (less style injection $\Rightarrow$ less content distortion). The ratio $\Delta\text{LPIPS}/\Delta\text{CLIP-S} \approx 8$ is therefore the ratio of the two loss curvatures on $\mathcal{M}_0$, not a property of any specific knob. Every capacity knob traces the same line because each perturbs only $\kappa$, and $\kappa$ is the only degree of freedom on $\mathcal{M}_0$.

\paragraph{What breaks the frontier.}
Only three structural degrees of freedom move the operating point off the 1:8 line: (i) the wavelet decomposition itself, which moves the model off $\mathcal{M}_0$ (Proposition~4); (ii) the solver order (Euler~$\to$~Heun, $+$0.0056 CLIP-S at fixed $\alpha$); (iii) the stochastic routing probability $p$. Capacity, blend, schedule, and training duration do not.

\section{Multi-Seed Stability}
\label{sec:multiseed}

\subsection{Three-Seed CLIP-S and LPIPS}
\label{sec:threeseed}

Table~\ref{tab:multiseed} reports SFM across three random seeds. The main paper reports the seed-1 value. The main paper states that multi-seed variance is below $0.002$ CLIP-S for SFM.

\begin{table}[h]
\centering
\caption{Multi-seed stability of SFM on the 750-pair protocol.}
\label{tab:multiseed}
\small
\begin{tabular}{@{}lcc@{}}
\toprule
Seed & CLIP-S $\uparrow$ & LPIPS $\downarrow$ \\
\midrule
1 (reported) & 0.7213 & 0.2868 \\
2 & TBR & TBR \\
3 & TBR & TBR \\
\midrule
Mean $\pm$ std & $0.7213 \pm \sigma_C$ & $0.2868 \pm \sigma_L$ \\
\bottomrule
\end{tabular}
% TODO: fill in seeds 2 and 3. Per the main paper, std_C < 0.002 CLIP-S and
% std_L < 0.005 LPIPS. Exact per-seed values to be reported upon code release;
% do not fabricate.
\end{table}

The standard deviation satisfies $\sigma_C < 0.002$ on CLIP-S and $\sigma_L < 0.005$ on LPIPS, consistent with the bound stated in the main paper's AAAI Reproducibility Checklist.

\subsection{Variance Explanation}
\label{sec:variance}

The low variance is structural. With LL locked, the flow-matching loss on the high-frequency subbands, $\mathcal{L}_{\text{HF}}(\theta) = \mathbb{E}_t \|v_\theta^{\text{HF}}(z_t, t, s) - u_t^{\text{HF}}\|^2$, is approximately convex near the optimum (Section~4.3 of the main paper). The target $u_t^{\text{HF}}$ carries at most $\epsilon$ of the content energy (Proposition~1), so the velocity field fits a low-complexity brushwork target. Different seeds therefore converge to nearby optima, yielding small metric dispersion. By contrast, the Euclidean baseline optimizes a non-convex reconciliation between content and style gradients on $\mathcal{M}_0$ and shows larger seed-to-seed spread.

\section{Additional Experimental Details}
\label{sec:exp_details}

\subsection{Dataset Details}
\label{sec:dataset}

We use the Distinct5-WikiArt-512 benchmark: five artistic domains --- Early Renaissance, Impressionism, Minimalism, Rococo, Ukiyo-e --- each with $3{,}600$ training images and $150$ test images. Inputs are SDXL VAE latents at $4 \times 32 \times 32$ (decoded to $512 \times 512$ images). Training uses $18{,}000$ latents; evaluation uses $750$ pairs ($5$ source domains $\times$ $5$ target domains $\times$ $30$ source images).

\subsection{Evaluation Protocol}
\label{sec:eval}

CLIP-S is the cosine similarity between CLIP ViT-B/32 embeddings of the output and a target-style prototype (higher is better). LPIPS is the AlexNet perceptual distance between the output and the content image (lower is better). We use the HuggingFace \texttt{transformers} CLIP backend and the standard AlexNet LPIPS. The Identity baseline (no transfer) gives CLIP-S $0.6933$ and LPIPS $0.0000$; it is both the content-preservation ceiling and the style-failure floor. All trainable baselines are evaluated under this unified 750-pair protocol.

\subsection{Baseline Training Details}
\label{sec:baselines}

All trainable baselines are trained to convergence on the same Distinct5 training set. Training times: CUT, $322.6$ min on an RTX 3090; SaMam, $436$ min ($20{,}000$ steps) on an RTX 3090; SaMST, $39.5$ min on an RTX 3090. SFM trains in $3.08$ min on a single RTX 3060 (12\,GB).

\section{Failure Cases and Limitations}
\label{sec:failures}

\subsection{Minimalism Target Failure}
\label{sec:minimalism}

Minimalism is largely low-frequency color blocking. Because base locking fixes LL (the global tone channel), SFM cannot shift the global palette strongly. To match the Minimalism target prototype, the model compensates by over-writing high-frequency structure, which inflates LPIPS to $0.5511$ (the highest of the five target domains). The trade-off is by design: SFM targets real-texture transfer, not global-palette replacement. A dedicated low-frequency style channel would be needed for palette-driven styles.

\subsection{High-Frequency Content Overwrite}
\label{sec:hf_overwrite}

Style is injected by per-subband WCT on LH, HL, and HH. When the content image itself carries strong high-frequency structure (text, fine line work, dense patterns), the injection overwrites that structure with the target-domain brushwork statistics. This is the converse of the content-preservation guarantee in Proposition~1: the guarantee holds for content concentrated in LL, not for content whose energy extends into LH/HL/HH. Mitigation would require a content-conditional blend per subband, which we leave to future work.

\subsection{Domain-Level Style, Not Per-Image Style}
\label{sec:domain_level}

SFM transfers the style \emph{distribution} of a domain, conditioned on a per-domain learnable token bank. It does not transfer the unique brushwork of a single reference painting. Two content images stylized to the same target domain therefore share a domain-level texture signature rather than an image-specific one. Per-image style would require a style encoder conditioned on a single reference image, which is outside the current architecture and is left to future work.

\end{document}
